\subsection{正则局部环的基本同调性质}

命题 1.- 设 A 为 Noether 正则局部环，n 为其维数。则有 \(\mathrm{dh(A)} = n\)，且对每个整数 \(i \geqslant 0\)，

\[
\left[ \operatorname{Ext} _ {A} ^ {i} \left(\kappa_ {A}, \kappa_ {A}\right): \kappa_ {A} \right] = \left[ \operatorname{Tor} _ {i} ^ {A} \left(\kappa_ {A}, \kappa_ {A}\right): \kappa_ {A} \right] = \binom {n} {i}.
\]

设 \(\mathbf{x} = (x_{1},\ldots ,x_{n})\) 为 A 的一个坐标系 (VIII, § 5, 第 1 节, 定义 1)。序列 \(\mathbf{x}\) 生成 \(\mathfrak{m}_{\mathrm{A}}\) 且对 A 是完全割的（同上引文, 第 2 节, 定理 1）。Koszul 复形 \(\mathbf{K}_{\bullet}(\mathbf{x},\mathrm{A})\) 是 \(\kappa_{\mathrm{A}}\) 的一个自由分解 (A, X, 第 159 页, 注 3)，其微分模 \(\mathfrak{m}_{\mathrm{A}}\) 为零。因此，对每个整数 \(i\geqslant 0\)，有（ \(\S 3\) , 第 3 节, 公式 (1)）

\[
\left[ \operatorname{Ext} _ {A} ^ {i} (\kappa_ {A}, \kappa_ {A}): \kappa_ {A} \right] = \left[ \operatorname{Tor} _ {i} ^ {A} (\kappa_ {A}, \kappa_ {A}): \kappa_ {A} \right] = \operatorname{rg} _ {A} (\mathbf {K} _ {i} (\mathbf {x}, A)) = \binom {n} {i}.
\]

于是由 § 3, 第 3 节, 命题 4 的推论 1 得出 dh(A) = n。

命题 2.--- Noether 正则局部环是因子分解整环。

由命题 1，Noether 正则局部环上的每个有限生成模都具有一个有限长度的投射分解，它由有限生成投射模组成，从而由自由模组成 (II, § 3, 第 2 节, 命题 5 的推论 2)。于是由 VII, § 4, 第 7 节, 命题 16 的推论 3，这样的环是因子分解整环。

命题 3.- 设 A 为 Noether 正则局部环，M 为非零的有限生成 A-模。其投射维数有限，且有

\[
\mathrm{dp} _ {\mathrm{A}} (\mathrm{M}) + \operatorname{prof} _ {\mathrm{A}} (\mathrm{M}) = \dim (\mathrm{A}).
\]

事实上，M 具有有限投射维数（命题 1），且有 \(\operatorname{prof}(A) = \dim(A)\)，因为 A 是 Macaulay 环 ( \(\S\) 2, 第 5 节, 例 7)。然后应用 \(\S\) 3, 第 5 节的定理 1。

推论 1. 有 \(\mathrm{dp}_{\mathrm{A}}(\mathrm{M}) \geqslant \dim(\mathrm{A}) - \dim(\mathrm{M})\) ；为使等号成立，必须且只需 M 是 Macaulay 的。

推论 2.--- 为使 A-模 M 是自由的，必须且只需它是 Macaulay 的且维数为 \(\dim(A)\)，或者等价地，它的深度 \(\geqslant \dim(A)\) 。

推论 3.--- 维数为 2 的 Noether 正则局部环上的每个有限生成自反模都是自由的。

事实上，Noether 正则局部环是整闭的 (VIII, § 5, 第 2 节, 定理 1 的推论 1)。于是推论 3 由推论 2 及 § 1, 第 10 节的命题 16 得出。

推论 4.--- 设 \(\rho: A \to B\) 为 Noether 局部环之间的局部同态。假设 A 正则且 \(\rho\) 使 B 成为有限生成 A-模。则有 \(\mathrm{dp}_{\mathrm{A}}(\mathrm{B}) \geqslant \dim(\mathrm{A}) - \dim(\mathrm{B})\) 。为使 B 是 Macaulay 环，必须且只需 \(\mathrm{dp}_{\mathrm{A}}(\mathrm{B}) = \dim(\mathrm{A}) - \dim(\mathrm{B})\) 。为使 B 是维数等于 \(\dim(\mathrm{A})\) 的 Macaulay 环，必须且只需 A-模 B 是自由的。

事实上，有 \(\dim(B) = \dim_A(B)\) (VIII, § 2, 第 3 节, 定理 1)；此外，B 是 Macaulay 环当且仅当它是 Macaulay A-模 (§ 2, 第 6 节, 命题 8)。因此只需应用推论 1 与推论 2。

注.--- 推论 4 使我们能在若干重要情形刻画 Macaulay 局部环。设 A 为 Noether 局部环；它是 Macaulay 环当且仅当其完备化 \(\widehat{\mathsf{A}}\) 是 Macaulay 环 (§ 2, 第 7 节, 命题 9 的推论 4)。以下设局部环 A 是完备的，并令 \(d = \dim(\mathsf{A})\) 。

\begin{enumerate}
\def\labelenumi{\alph{enumi})}
\tightlist
\item
  设 A 具有一个子域。于是它具有一个系数子域 K (IX, § 3, 第 3 节)，且存在一个形式幂级数代数 E = K{[}{[}T₁, \ldots, Tₙ{]}{]} 以及一个 K-代数的满射同态 E → A（同上引文）；还存在一个形式幂级数代数 E' = K{[}{[}T₁, \ldots, T\_d{]}{]} 以及一个 K-代数的单射局部同态 E' → A，使 A 是 E' 上的有限代数（同上引文）。下列性质等价：
  \begin{enumerate}
  \def\labelenumii{(\roman{enumii})}
  \item
    A 是 Macaulay 环；
  \item
    有 \(\mathrm{dp}_{\mathrm{E}}(\mathrm{A}) = n - d\)；
  \item
    A 是自由 E'-模。
  \end{enumerate}

\item
  设 A 的剩余域特征为 p \textgreater{} 0。存在一个长度为 +∞ 的 p-环，其剩余域为 \(\kappa_{\Lambda}\) (IX, § 2, 第 3 节, 命题 5)。设 C 为这样一个环；存在一个形式幂级数代数 \(E = C[[T_{1}, \ldots, T_{n}]]\) 以及一个满射同态 \(\rho : E \to \Lambda\) (IX, § 2, 第 5 节, 定理 3)。下列性质等价：
  \begin{enumerate}
  \def\labelenumii{(\roman{enumii})}
  \item
    A 是 Macaulay 环；
  \item
    有 \(\mathrm{dp}_{\mathbf{E}}(\mathrm{A}) = n + 1 - d\)；
  \end{enumerate}

再设 \(p1_{\Lambda}\) 在 A 中不是零因子；于是存在一个形式幂级数代数 \(E' = K[[T_1, \ldots, T_{d-1}]]\) 以及一个 K-代数的单射局部同态 \(E' \to A\)，使 A 是 \(E'\) 上的有限代数（同上引文）。局部环 \(E'\) 是正则的，维数为 \(n+1\) (VIII, § 5, 第 5 节, 例 2)。上述条件也等价于

  \begin{enumerate}
  \def\labelenumii{(\roman{enumii})}
  \setcounter{enumii}{2}
  \tightlist
  \item
    A 是自由 E'-模。
  \end{enumerate}
\end{enumerate}

关于模的情形的类似结果，参见 § 5, 第 5 节。

